\section{Related work}
\label{sec:related}

We consider a world model to be a generative system that autoregressively generates image or video frames, conditioned on actions or camera pose. This work builds on several important concepts in the generative AI literature, which we briefly review in the following.  

\paragraph{Image and video generation.}
% Diffusion models and flow matching have been established as state-of-the-art methods for image and video generation [CITE a bunch of stuff]. 

In recent years, diffusion models have emerged as the state-of-the-art paradigm for image generation, surpassing the performance of GAN-based methods in both fidelity and diversity by modeling complex distributions through iterative denoising~\cite{Lee2018StochasticAV, Mathieu2015DeepMV, Tulyakov2017MoCoGANDM, Vondrick2017GeneratingTF, Dhariwal2021DiffusionMB}. These successes have naturally extended to the task of video generation, as diffusion-based architectures have been adapted to include the temporal domain, enabling the generation of high-quality video clips spanning the order of tens to hundreds of frames~\cite{Zheng2024OpenSoraDE,Lin2024OpenSoraPO,Yang2024CogVideoXTD,Qiu2025SkyReelsA1EP,Fei2025SkyReelsA2CA,Chen2025SkyReelsV2IF,Blattmann2023StableVD}. Additional advancements in efficiency and controllability have also followed, through techniques such as improved tokenization~\cite{Tang2024VidTokAV,Xie2024SANAEH} and flow-matching objectives~\cite{Lipman2022FlowMF,Jin2024PyramidalFM}.

\paragraph{Autoregressive video generation.}
% goal: long generation; causal attention, diffusion forcing, sliding window approach. potentially briefly talk about relative benefits and disadvantages of these methods (no need to contrast with ours).

To support generation of longer videos with online inference capabilities, autoregressive (AR) approaches and architectures have been introduced in the diffusion framework~\cite{sun2024autoregressive,kondratyuk2024videopoet,yan2021videogpt,wang2024loong}. The autoregressive regime is also particularly intuitive, given the natural temporal ordering of video data. Taking inspiration from LLMs, state-of-the-art methods model video data by processing frames into spatio-temporal tokens~\cite{Lin2024OpenSoraPO,Zheng2024OpenSoraDE,Yang2024CogVideoXTD,Tang2024VidTokAV,Fei2025SkyReelsA2CA,Qiu2025SkyReelsA1EP,Chen2025SkyReelsV2IF}, learning to autoregressively generate new tokens until a full video is formed. This is in stark contrast with conventional diffusion-based frameworks, which choose to iteratively denoise an entire image/video at once~\cite{blattmann2023stable,li2022efficient,ho2022video,singer2022make,yang2024cogvideox,kong2024hunyuanvideo,hacohen2024ltx, oshima2024ssm}. Recent approaches include training conditional diffusion models that generate new frames given a set of previous clean frames, allowing AR generation simply by passing newly generated frames as context for future frames~\cite{ho2022video,he2022latent,chen2023seine,jin2024pyramidal,zhang2024extdm,gao2024ca2,gao2024vid}. Others modify the diffusion objective by assigning independent noise levels for each frame during training, which also supports AR inference~\cite{chen2024diffusion,yin2024causvid,song2025history}. In practice, these methods can be used for infinite-length generations by utilizing a sliding window context, but suffers from limited memory and drift.


\paragraph{Controlled video generation.} 
Controlled or conditional video generation is a core component of world simulators. 
A series of works have explored explicit camera control, including~\cite{wang2024motionctrl,he2025cameractrl,he2025cameractrl2,he2025cameractrl2,bahmani2025vd3d,bahmani2025ac3d}, which enable novel-view or multi-view generation by disentangling camera trajectories from dynamic content. Most of these approaches take either a single image or a text prompt as input, and often struggle to maintain long-term consistency, especially in dynamic scenes. 
More recent efforts~\cite{zhang2024recapture,bai2025syncammaster,bai2025recammaster} focus on re-filming or generating synchronized views in dynamic settings, further advancing controllability. 
Beyond direct camera injection, structural conditioning via point clouds, tracking, or 3D-aware priors has proven effective for improving spatial consistency and trajectory alignment~\cite{yu2024viewcrafter,ren2025gen3c,gu2025das,yu2025trajectorycrafter, chen2025flexworld}.
In addition to spatial and motion-level control, some models support action-based or scene-level conditioning~\cite{parkerholder2024genie2,oasis2024}, where either structured action vocabularies or segmented text prompts serve as high-level drivers of video progression.

\paragraph{Long-context video generation.}
While advances in autoregressive methods have allowed for generation of longer, and even potentially inifinite-length videos~\cite{Chen2025SkyReelsV2IF}, the reliance on sliding context windows limit their effective memory. The straightforward approach of increasing memory by training with longer context lengths is effective but computationally demanding. Recent works have explored efficient architectural alternatives to attention-based transformers, such as state-space models and linear attention~\cite{Wang2024LinGenTH,Liu2024LinFusion1G}. While efficient, such methods often perform worse than conventional diffusion transformer architectures. Other concurrent works have proposed compressing prior frames to lower context length~\cite{Zhang2025PackingIF,Gu2025LongContextAV}, subsequently lowering computational demand at the cost of information lost in context frames. A separate branch of works ground previously generated frames in a 3D representation (e.g., point clouds ~\cite{yu2024viewcrafter,gu2025das,yu2025trajectorycrafter,ren2025gen3c}) and query this representation based on the camera position of future frames. These methods allow for precise camera control, but struggle to handle dynamic scenes with complex motions.

In this work, we adopt the conventional approach of concatenating recent context frames as working memory in diffusion-based generation to autoregressively generate multiple future frames. Meanwhile, we iteratively predict and filter the static point maps of newly generated frames to update a global spatial memory, which serves as geometric guidance for long-term generation. We further incorporate a sparse set of historical reference frames as an episodic memory.

% 1-2 sentences on how our model combines existing components from the above (i.e., which image/video base model; what type of autoregressive-ness; what type of control or action conditioning. optionally say that our core contribution is the memory part, which is (in principle) compatible with any world model.

% Brief 1-3 sentence summary of how our memory approach is similar or different (definitely better) than existing approaches. 


\endinput

% \paragraph{Video World Models.}
% % \tong{Need to first narrow down to the video generative models to rule out the world model branch in autonomous driving}
% World models aim to predict future states when provided a set of previous states and corresponding actions, enabling applications in robotic planning, simulation, and interactive experiences. Specifically, video world models represents this state through a set of video frames.
% %
% While video world models have shown promise in earlier works~\cite{ha2018recurrent,menapace2021playable,chiappa2017recurrent,oh2015action}, they often rely on architectures such as RNNs~\cite{hochreiter1997long} and variational autoencoders~\cite{kingma2013auto}. Leaning on a more scalable approach, recent works have moved towards transformer-based diffusion models, resulting in impressive video generation quality.
% % 
% Further works have seen advancement in several key directions: training models without relying on ground truth action data~\cite{Bruce2024GenieGI,menapace2021playable}, scaling to more complex environments and control systems~\cite{valevski2024diffusion,yang2023learning,xiang2024pandora,hu2023gaia}, achieving real-time inference~\cite{feng2024matrix,oasis2024}, and generating open-world environments~\cite{yu2025gamefactory,che2024gamegen,feng2024matrix}.
% %
% To enable video generation over long horizons, these models utilize transformers trained on shorter context lengths in a sliding-window manner. Given the limited attention window of such models, generated frames beyond this window are discarded. This inherently limits the models ability to remember past context beyond the window, which usually only spans several frames.
% %
% On top of paying attention to past frames, our method stores a grounded 3D representation of the generated world through an auto-regressively updated point map. Rather than relying solely on past frames, future frames can access relevant context by back-projecting this structured memory into the current camera view. This design allows our model to retain a long effective memory by storing prior information in an explicit state (3D points), overcoming the limitations of having a fixed attention window.

% \noindent{\textbf{Generative Novel View Synthesis}}
% Generative novel view synthesis has seen rapid progress, particularly with diffusion-based models capable of controlling camera motion. 
% %
% Early approaches like 4DiM~\cite{watson20254DiM} enabled novel view generation by modeling space-time jointly, but were limited to per-frame synthesis and struggled with dynamic content. Transformer-based video diffusion models such as VD3D~\cite{bahmani2025vd3d} and AC3D~\cite{bahmani2025ac3d} improved camera control by disentangling camera motion from scene dynamics, with AC3D introducing refined pose conditioning and training strategies to reduce drift. CameraCtrl-II~\cite{he2025cameractrl2} and Collaborative Video Diffusion~\cite{} incorporated pose-aware conditioning to enable multi-view video generation, though they lacked persistent scene memory. A key advancement came from TrajectoryCrafter~\cite{yu2025trajectorycrafter}, which introduced an explicit 3D memory via point clouds to guide novel-view rendering of monocular videos along user-defined paths, preserving content over time. Similarly, ReCamMaster~\cite{bai2025recammaster} and SynCamMaster~\cite{} used enhanced conditioning and synthetic training data to support consistent re-filming of dynamic scenes. 
% %
% However, these methods often struggle with complex dynamics or maintaining consistency when revisiting past views. In contrast, our method models explicit memory, combining long-term static memory with short-term dynamics and recent actions. This enables better 3D consistency and broader applicability to world modeling tasks beyond simple point-cloud conditioning.


% \noindent{\textbf{Learning-Based 3D Reconstruction}}
% Traditional structure-from-motion (SfM) pipelines such as COLMAP~\cite{schoenberger2016sfm} have long served as the gold standard for image-based 3D reconstruction due to their high accuracy, but they rely on incremental feature matching and bundle adjustment that do not scale well to large-scale or real-time applications. To overcome these limitations, learning-based approaches have emerged: the seminal DUT3R~\cite{wang2024dust3r} method established a new paradigm of dense matching and optimization in an end-to-end deep learning framework, eliminating the need for separate SfM or multi-view stereo modules. Building on this, recent approaches~\cite{zhang2025flare,wang2025vggt} introduce feed-forward architectures that significantly accelerate inference, achieving near real-time reconstruction without iterative optimization. To further process dynamic scenes, more recent advances~\cite{zhang2025monstr,li2024megasam,wang2025cut3r} further extend learning-based 3D reconstruction to dynamic scenes by incorporating recurrent models or persistent state, which improves efficiency and robustness in handling moving objects and changing environments. These advances make it feasible to perform online estimation of camera poses and point clouds during video capture or generation. Building on Mega-SAM, our method integrates an iterative video generation process with joint point cloud estimation and a dynamic-static disentanglement mechanism to preserve long-term static scene memory.

% Table generated by Excel2LaTeX from sheet 'Sheet1'
\begin{table}[t]
  \centering
  \caption{\textbf{Comparison with different spatial-aware video generative models}}
    \small
    \begin{tabular}{l|l|cccc}
    \toprule
    \multicolumn{2}{l|}{Methods} & Cam-Ctrl & Long-Mem & Action-Pred & S/D-Decouple \\
    \midrule
    Point condition & \cite{yu2024viewcrafter,gu2025das,yu2025trajectorycrafter,ren2025gen3c} & \ding{51} & \ding{55} & \ding{55} & \ding{55} \\
    \multirow{2}[0]{*}{Camera condition} & \cite{wang2024motionctrl,bahmani2025vd3d,bahmani2025ac3d,he2025cameractrl2} & \ding{51} & \ding{55} & \ding{51} & \ding{55} \\
          & \cite{zhang2024recapture,bai2025syncammaster,bai2025recammaster} & \ding{51} & \ding{55} & \ding{55} & \ding{55} \\
    Hidden state & \misscite{Po's}  & \ding{51} & \ding{51} & \ding{51} & \ding{55} \\
    History frames & \cite{xiao2025worldmem} & \ding{51} & \ding{51} & \ding{51} & \ding{55}  \\
    \midrule
    \multicolumn{2}{l|}{\textbf{Ours}} & \ding{51} & \ding{51} & \ding{51} & \ding{51} \\
    \bottomrule
    \end{tabular}%
  \label{tab:taxonomy}%
\end{table}%

% % Table generated by Excel2LaTeX from sheet 'Sheet1'
% \begin{table}[t]
%   \centering
%   \caption{\textbf{Comparison with different spatial-aware video generative models}}
%     \small
%     \begin{tabular}{l|l|ccccc}
%     \toprule
%     \multicolumn{2}{l|}{Methods} & Cam-Ctrl & Long-Mem & Action-Pred & S/D-Decouple & Domain \\
%     \midrule
%     Point condition & \cite{yu2024viewcrafter,gu2025das,yu2025trajectorycrafter,ren2025gen3c} & \ding{51} & \ding{55} & \ding{55} & \ding{55} & Real-world \\
%     \multirow{2}[0]{*}{Camera condition} & \cite{wang2024motionctrl,bahmani2025vd3d,bahmani2025ac3d,he2025cameractrl2} & \ding{51} & \ding{55} & \ding{51} & \ding{55} & Real-world \\
%           & \cite{zhang2024recapture,bai2025syncammaster,bai2025recammaster} & \ding{51} & \ding{55} & \ding{55} & \ding{55} & Real-world \\
%     Hidden state & \misscite{Po's}  & \ding{51} & \ding{51} & \ding{51} & \ding{55} & MineCraft \\
%     History frames & \cite{xiao2025worldmem} & \ding{51} & \ding{51} & \ding{51} & \ding{55} & MineCraft/Static scenes \\
%     \midrule
%     \multicolumn{2}{l|}{\textbf{Ours}} & \ding{51} & \ding{51} & \ding{51} & \ding{51} & Real world \\
%     \bottomrule
%     \end{tabular}%
%   \label{tab:taxonomy}%
% \end{table}%


% Recent advances in generative novel view synthesis have enabled remarkable control of camera trajectories in both static and dynamic scenes. Early 3D-aware generative models for static scenes injected geometric priors (e.g., latent 3D feature grids or NeRF-like representations) into diffusion models to ensure view consistency from limited inputs, but they assumed no scene motion. Building on this, 4DiM introduced a space-time diffusion framework trained on mixed 3D/4D data, pioneering novel view synthesis with temporal dynamics from as few as one image, albeit generating one frame at a time. Parallel efforts in text-to-video generation integrated camera pose conditioning into large diffusion models: initial works like CameraCtrl and MotionCtrl added plug-and-play pose encoders to pre-trained video Diffusion U-Nets, while VD3D and the subsequent AC3D leveraged transformer-based video diffusion backbones for finer 3D camera control. AC3D, in particular, analyzed the implicit camera understanding in video diffusion models to inject pose information only in low-frequency denoising layers and introduced training data of dynamic scenes with stationary cameras to disentangle camera motion from scene motion, achieving more precise camera moves without degrading visual quality. On the other hand, several methods explicitly incorporate memory or 3D scene representations to maintain consistency. For example, TrajectoryCrafter employs a dual-stream diffusion model that conditions on both the source video and an extracted point cloud of the scene – this explicit geometric memory allows precise re-filming of a monocular video along new user-defined trajectories, keeping the content coherent over time. Similarly, ReCamMaster re-renders in-the-wild videos from novel camera paths by harnessing a pre-trained text-to-video diffuser with a robust video-conditioning mechanism, boosted by training on a synthetic multi-camera video dataset to preserve multi-frame appearance and dynamic synchronization. Beyond single video re-rendering, multi-view video generation has been tackled by Collaborative Video Diffusion and SynCamMaster, which coordinate multiple diffusion models or a unified model to produce synchronized videos from diverse viewpoints of the same scene. These systems enforce cross-view consistency—often via a shared latent scene representation—so that dynamic objects and background details remain temporally and geometrically aligned across cameras. In summary, static-scene approaches achieve photorealistic novel views through implicit 3D reasoning but cannot handle moving content, whereas dynamic-scene methods introduce specialized conditioning, data augmentation, or explicit memory modules to capture independent object motion during camera movement. Our method differs from all prior works by explicitly targeting fully dynamic scenes and incorporating a dedicated memory mechanism, enabling it to maintain coherent scene appearance under complex camera trajectories where previous methods struggle.